\chapter{An implementation contract for the interpretation workbench}
\label{ch:implementation}

\section{Extend the existing product at a defined boundary}

The current repository already has source intake, RuleIR validation and execution,
event replay, Java generation, review records, build and release manifests, and a
local service interface. Its drafting provider is a deterministic fixture. The
new work is an interpretation investigation service between source intake and
approval of a candidate bundle. Replacing the entire runtime would discard useful
engineering evidence and enlarge the task unnecessarily.

The extension accepts a retained source packet, an activity profile and an
explicit investigation policy. It creates inventories, candidates, arguments,
discrepancies and bounded actions. Its output is a report and, when appropriate,
one or more validated draft bundles. The existing review and release services
remain responsible for approval and export, with an additional requirement that
the interpretation report and its unresolved issues be checked.

This boundary has two advantages. The interpretation service can evolve its
candidate-generation methods without changing deterministic production execution.
The Java path can be verified independently of any particular model provider.
A candidate that requires an unsupported operator is returned as an engineering
question rather than forcing a runtime change through an unreviewed model output.

The companion files under \srcpath{docs/monograph/contracts} specify a proposed
\texttt{interpretation.v1} extension. They are design contracts and examples, not
a claim that the service is already implemented. The ordered task file under
\srcpath{docs/monograph} identifies dependencies and acceptance conditions. An
implementation agent should read this chapter with those files and the existing
v0.1 runtime contracts; it should not infer additional authority from illustrative
code in the manuscript.

\section{Separate durable records from generated prose}

The service needs a small set of durable record types. A run binds the source
packet, activity profile, policy and required ensemble roles. A candidate binds a
proposed meaning to a formal representation and its assumptions. An issue records
a discrepancy and its evidence need. An action records one authorized external
operation. A report reconciles the run's evidence and terminal status. Review
decisions and source coverage records connect these objects to human authority
and document completeness.

Generated prose is a field within these records, not their controlling structure.
A model may write an explanation saying an issue is resolved, but the service
changes the issue's resolution status only through a validated transition with
the required evidence. A sentence saying ``approved by compliance'' has no effect
unless an authenticated review record exists for the exact object.

Identifiers are stable opaque strings within a namespace; content hashes identify
immutable payloads. These roles should not be confused. A run identifier remains
stable while events accumulate, whereas each immutable candidate revision has a
new content identity. A human-readable title can change without becoming the
primary key. References use identifiers and, where integrity requires it, the
expected content hash.

Timestamps use the project's canonical UTC format. Durations are integer
milliseconds or seconds with explicitly named units. Counts are nonnegative
integers. Monetary reservations use exact minor units and a named billing
currency. No field called \texttt{cost} should leave the unit or estimated-versus-
actual distinction implicit. This discipline prevents the controller from
comparing incompatible budgets.

The schemas reject unknown properties. This catches misspelled limits and
unrecognized status fields rather than ignoring them. Structural validation is
only the first step: the service also verifies references, state transitions,
source compatibility, budget arithmetic and role authority. JSON Schema cannot
prove that a source passage supports an interpretation or that two referenced
objects belong to the same run.

\section{The run and policy records}

A run records its immutable configuration hash, source packet hash, profile
identifier, required roles, creation time, deadline and current lifecycle state.
It also references its candidate, issue, action and report collections. Mutable
projections such as counters are derived from durable events or updated
transactionally with them; they are never trusted solely because a client submits
a number.

The lifecycle begins with initialization, proceeds through initial proposals and
resolution rounds, and ends in ready-for-review, blocked-unresolved, cancelled or
failed-integrity. A ready-for-review run is complete enough to present, not
released. A run that is blocked can still contain useful validated candidates.
The API returns them with the block, rather than suppressing all information
because the overall investigation is incomplete.

The policy supplies all material resource and coverage choices. It names required
roles, maximum initial actions, maximum resolution rounds, per-root and global
action caps, no-progress limit, action timeout, run duration, candidate count,
argument-node count and dependency depth. It also declares whether monetary and
token caps are enforced. Configuration validation checks internal consistency,
such as requiring the initial-action cap not to exceed the global cap.

The reference fixture policy is deliberately explicit:
\begin{lstlisting}[caption={Reference limits for deterministic controller fixtures.}]
{
  "schema_version": "interpretation.v1",
  "policy_id": "fixture-policy.v1",
  "use": "DETERMINISTIC_FIXTURE_ONLY",
  "max_initial_actions": 4,
  "max_resolution_rounds": 3,
  "max_actions_per_issue_root": 3,
  "max_actions_total": 12,
  "max_no_progress_rounds": 2,
  "action_timeout_seconds": 30,
  "run_deadline_seconds": 600,
  "max_candidates": 24,
  "max_argument_nodes": 200,
  "max_dependency_depth": 6,
  "token_cap": null,
  "cost_cap_minor_units": null
}
\end{lstlisting}

The candidate, graph and dependency bounds make local processing tractable for
fixtures. Reaching one opens or preserves an incompleteness issue. These numbers
are convenience choices for engineering tests, not selected production defaults.
A paid provider configuration must add enforceable cost and token reservations,
and its deployment policy must explain why the chosen limits are sufficient for
the intended workload.

A policy version is immutable once used. Changing a cap during a run requires an
explicit authorized amendment event if that capability is later implemented;
the first prototype should instead require a successor run. This keeps budget
semantics simple and prevents a worker from extending its own authority. The
successor links to the prior report and explains why additional work is justified.

\section{The candidate record carries the interpretive commitment}

A candidate has a parent revision, source references, subject unit, controlled-
language statement, formal bundle reference, assumptions and interpretation
family. The family identifies a consequential choice, such as component versus
whole-offer assessment. It is not merely a cluster label assigned by a text
embedding. A candidate can belong to several dimensions, but its changed choices
must be explicit relative to its parent.

Each assumption has a statement, provenance and status. Possible statuses include
supported by inspected source, factual evidence required, provisional modeling
choice, disputed and rejected. A known fact in a runtime snapshot must not be
created directly from a disputed assumption. The candidate-to-fact mapping is a
reviewed transformation with its own evidence requirements.

Supporting and opposing arguments reference source spans or other arguments.
An argument distinguishes premises, inference rule and conclusion, and labels
strict versus defeasible inference where that distinction is used. If an
argumentation engine is not yet implemented, these records can still support a
structured review interface. The service must not label the resulting graph as
formally evaluated until the selected semantics has actually been run.

Candidate status distinguishes active, deferred, defeated, unsupported and
selected-for-review. Deferred means not yet sufficiently investigated, not false.
Defeated requires a recorded reason under the chosen argument and authority
policy. Unsupported means the current record lacks adequate support. Selected
for review is a presentation decision; only a separate authenticated decision can
approve the candidate's meaning.

The record includes score components with definitions and calibration status.
The initial contract permits a scheduling priority and categorical evidence
states. It explicitly leaves probability of legal correctness null. Later
calibration can add a versioned, narrowly defined score, but cannot reinterpret
old uncalibrated priorities as probabilities after the fact.

Source and assumption identity survive formal deduplication. If two candidates
share an identical RuleIR body, the body can be stored once. Their distinct
assumptions and source mappings remain separate candidate records. This prevents
a compiler optimization from erasing the very disagreement the workbench is
intended to expose.

\section{The issue and action records make automatic work accountable}

An issue names its root, parent if any, run, source units, affected candidates,
semantic dimension and materiality. It records the observed discrepancy and the
question that would resolve it. Processing state and resolution state are separate
fields. A terminal processing state can coexist with unresolved meaning.

Resolution evidence is typed. A source-based resolution references retained spans
and the reviewer or deterministic rule that accepted them. An adjudication
references the authenticated decision. An equivalence resolution references the
comparison report, domain and confirmation that no distinct source commitment
remains. A bare model statement is not a valid resolution-evidence type.

An action names the run, issue root or initial role, round, permitted action type,
input hashes, adapter, timeout, reservation and current state. It also stores the
idempotency key, lease owner, fencing token and terminal result reference. The
action's question and success criterion are immutable after dispatch. Changing
them creates another counted action.

The action result separates transport success from substantive progress. A model
request can return valid JSON yet fail to provide any supported new proposition.
A retrieval can succeed technically while returning the wrong version. A solver
can complete normally with an unknown result. The adapter reports its technical
outcome, and the controller validates whether the result answers the action's
question.

External text is untrusted data. It cannot name a new tool and cause execution,
change a budget, alter role permissions or select an arbitrary filesystem path.
Permitted actions come from a server-side registry. The model can propose an
action from that registry, but the controller validates its applicability and
reserves resources before dispatch. This is both a security boundary and a way to
keep the method's stopping proof connected to actual behavior.

The issue record retains attempted actions even when they fail. It is often useful
to know that an official dependency was unavailable at the assessment cut-off,
that two classification prompts produced no new evidence, or that a formal check
was outside the supported fragment. Those facts explain why the investigation
stopped and what would justify reopening it.

\section{Coverage records connect source units to consequences}

A coverage record maps each inventoried unit to one or more dispositions:
implemented rule, contextual explanation, definition dependency, separate
obligation, reviewed out-of-scope provision, or unresolved question. The record
includes the reason and responsible reviewer where a legal judgment is involved.
A paragraph can have more than one role; it may contain both a definition and an
operative condition.

Coverage validation first checks that the record refers to the same source hash
as the inventory. It then checks that every inventoried unit appears, every
referenced rule or issue exists, and any claim of implemented coverage has a
validated bundle and verification evidence. It cannot decide the substantive
adequacy of a legal disposition from structure alone, so that judgment remains
an explicit review obligation.

Dependencies are edges with types and versions. An incorporated definition is
stronger than a background analogy in the sense that missing it can directly
prevent interpretation of the selected rule. The impact service uses these edges
to identify affected candidates when a source changes. A generic document list
would not tell the service which conclusions depend on which definitions.

The source inventory itself is versioned and reviewed. An improved extractor may
find a previously missed footnote without any change to the official bytes. That
creates a new derivative and inventory revision, not a new official source
publication. Candidates whose coverage was assessed against the older inventory
need reconsideration. Distinguishing source revision from extraction revision
makes that event explainable.

For the second circular, the executable scope remains paragraph 10. The coverage
record must retain all eighteen paragraphs and four footnotes with their
respective dispositions. A machine-readable count of one implemented paragraph
and seventeen other paragraph dispositions is more honest than a banner saying
``circular translated'' without a denominator. The report can still present the
selected slice as a successful engineering example while showing why whole-
document release is not established.

\section{Storage transactions and event history}

The first implementation can extend the existing transactional store. Tables for
runs, candidates, issues, actions, coverage revisions, reports and review decisions
store canonical payloads with indexed identity and relationship fields. Foreign
keys protect references, unique constraints protect idempotency, and transactions
protect coupled state changes. The exact storage engine can remain the existing
local choice for the prototype; bank deployment requires its own availability and
administrative-control assessment.

The most important transaction reserves an action. It reads the current run and
root counters under a lock or equivalent serializable mechanism, checks all
limits, inserts the action and outbox row, increments counters and appends the
audit event. If any step fails, none commits. A later worker cannot dispatch an
action that lacks the committed reservation.

A resolution transaction similarly checks the issue revision, validates the
resolution evidence, updates the issue projection and invalidates dependent
results when necessary. It does not modify immutable candidate or source blobs.
The event history records what changed and why, while the current projection
supports efficient queries. Recovery can rebuild projections or verify them
against the event sequence.

Client idempotency is scoped by authenticated caller and operation. Repeating the
same key with the same canonical payload returns the existing result. Reusing the
key with a different payload returns conflict. The service must not accept a
client-supplied role in the payload as authority. Authentication determines the
caller and its permitted operations.

The local audit log is useful for reproducibility but is not automatically
 tamper-resistant against a database administrator. A bank deployment may require
external append-only storage, signed checkpoints or equivalent controls. Those
choices are part of the operational security design, not properties conferred by
calling a table an audit log. The prototype should preserve enough identity and
ordering information to integrate such controls later.

\section{An API with reviewable state transitions}

The proposed HTTP extension uses asynchronous jobs for investigation and explicit
resources for review. Creating a run requires the source packet, profile,
configuration and idempotency key. The server validates all references and returns
a run identifier and job reference. It does not block the request until every
model has completed.

\begin{longtable}{@{}P{49mm}P{41mm}P{58mm}@{}}
\caption{Proposed interpretation-service operations. Existing v0.1 endpoints remain separate.}\label{tab:interpretation-api}\\
\toprule
Operation & Purpose & Required behavior\\
\midrule\endfirsthead
\toprule Operation & Purpose & Required behavior\\\midrule\endhead
\path|POST /interpretation-runs| & Create a bounded investigation & Validate source/profile/policy; authenticate caller; idempotent creation\\
\path|GET /interpretation-runs/{id}| & Inspect current state & Return counters, phase, terminal reason and linked resources\\
\path|POST /interpretation-runs/{id}/cancel| & Stop new automatic work & Preserve outstanding actions and produce a cancelled report\\
\path|GET /interpretation-runs/{id}/report| & Retrieve a report version & Expose blocked and failed reports as well as review-ready ones\\
\path|GET /interpretation-issues/{id}| & Inspect a discrepancy & Return candidates, evidence need, attempts and status separately\\
\path|POST /interpretation-issues/{id}/decisions| & Record human adjudication & Require reviewer authority, expected revision and exact evidence links\\
\path|POST /interpretation-runs/{id}/successors| & Authorize further work & Require reason, new policy and lineage; never reset the old run\\
\bottomrule
\end{longtable}

These paths are proposed additions, not descriptions of endpoints already present
in the retained OpenAPI file. The implementation task must generate a new OpenAPI
contract from actual request models and test it against the service. The book's
path table states responsibilities; the machine-readable schema files state the
record shapes.

Malformed requests return boundary errors without creating partial runs. Missing
references and incompatible source versions return semantic validation errors.
A stale expected revision returns conflict so the reviewer can reload the changed
issue. Unauthorized decisions fail without writing a review or audit entry that
suggests acceptance. The service records denied-operation security evidence
through its appropriate logging channel without pretending that the mutation
succeeded.

Pagination must preserve completeness. Collection responses include stable cursors
and total or explicitly unknown counts; reports bind the complete underlying
collection hashes. A user-interface page showing twenty issues cannot imply that
only twenty exist. Exports include all issue records or a manifest identifying
every page, with integrity checks that detect omission.

\section{The review interface should expose the decisive difference}

A reviewer should begin with the source passage and the decision under review,
not a wall of model transcripts. The interface shows the activity and subject
profile, the candidate's controlled-language rule, and the facts needed to apply
it. A neighboring comparison shows the strongest surviving alternative and a
scenario where the outcomes differ.

The discrepancy view explains the type of uncertainty. Missing factual evidence
should lead to a data request; unresolved meaning should lead to a legal question;
unsupported runtime operations should lead to engineering work. Mixing these
under a single confidence badge makes review slower and encourages the wrong
remedy. The interface can use concise status labels, with the full evidence one
step away.

Source references open the retained version and highlight the exact span. A live
link is useful for checking current authority, but it should not replace the
retained version used in the run. If the live document differs, the interface
shows the change and offers a successor investigation. This prevents a reviewer
from unknowingly approving a candidate against a document that changed after
its creation.

A decision form asks for the proposition being accepted, the evidence relied on,
the treatment of the strongest objection, the scope of the decision and any
conditions. It does not ask only for a confidence slider. The authenticated role
and exact candidate/report hashes are supplied by the service. The reviewer
cannot approve an unseen revision through a stale browser tab.

Unresolved reports remain useful work products. They can present a focused
question for the next authorized reviewer, show which offers are affected and
identify a temporary operational policy proposal. The interface should make clear
whether that policy is a bank choice or a conclusion about the regulator's rule.
This distinction is central to keeping cautious operations from becoming false
statements of law.

\section{Provider adapters are deliberately narrow}

A provider adapter accepts a typed task and returns a bounded response plus usage
and execution metadata. It has no direct database credentials for approvals,
release records or budgets. It cannot invoke arbitrary tools selected in model
text. The controller supplies the exact retained sources and permitted retrieval
scope, and the adapter reports which inputs were actually sent.

The initial interface can support three adapter kinds: scripted fixtures,
retrieval tools and model proposals. Scripted fixtures are essential because they
make failures reproducible. Retrieval adapters return bytes and acquisition
metadata for validation. Model adapters return structured proposals and source
references that the controller checks. All adapters use the same action identity
and timeout accounting.

The adapter must distinguish refusal, abstention, malformed output, transport
failure, timeout and unknown remote result. Treating them all as an empty answer
would lose the reason a mandatory role is incomplete. It must also report hidden
provider behavior, such as automatic retries or server-side tool execution, when
that behavior affects resource accounting or evidence provenance. An integration
that cannot bound these behaviors is unsuitable for the unattended profile.

Prompt templates are versioned source files. They state the role, source
perimeter, required structured fields and prohibition on treating source text as
instructions to the system. The prompt should ask for assumptions and opposing
reasons, but those requests do not make the response reliable by themselves.
Validation and independent evidence remain necessary.

The service stores provider and model identifiers, relevant sampling settings,
template version and source hashes. A provider alias that changes its underlying
model must be recorded as such. Reproducibility may be limited even with these
fields, so reports distinguish replay of a retained response from regeneration of
a fresh model response. Deterministic execution of an accepted bundle does not
depend on reproducing the model's original generation.

\section{Security protects the meaning of the evidence}

The source packet can contain malicious or misleading instructions, whether
through a compromised document, an uploaded file or quoted text. The interpreter
must treat them as source content, not commands. A paragraph saying to ignore
prior rules cannot change the controller's budget or authentication policy. The
model's task is to analyze that text within the source context; tool permissions
come from the server.

Acquisition should retain the existing restrictions on HTTPS destinations,
redirects, public addresses, response size and time. A user-supplied URL does not
establish official authority. Uploaded documents remain unverified until a
separate process establishes provenance. A source that is useful for comparison
may still be ineligible to support release.

PDF parsing needs an isolation boundary in a bank deployment. The current local
prototype's bounded parser is not an operating-system sandbox for arbitrary
hostile files. A production intake worker should have limited filesystem and
network access, resource limits and a narrow output contract. This protects the
workbench without changing the legal interpretation method.

Bank factual data and public regulatory sources have different access needs.
The initial ensemble prototype should use public sources and synthetic scenarios.
Sending client or transaction data to a model provider requires the bank's
approved data-handling arrangement. The architecture supports minimizing that
exposure: legal interpretation can often use abstract fact patterns, while
production evaluation remains inside the bank with the generated Java package.

Logs must avoid becoming an uncontrolled copy of sensitive evidence. Store
references and necessary metadata in operational logs, retain full evidence in
access-controlled stores, and make redaction policies explicit. A redacted report
must preserve enough structure to explain its limitations. Redaction cannot
silently remove the only evidence supporting an approval while leaving the
approval summary apparently complete.

Security tests should use concrete attacks: source text that requests a tool
call, candidate output containing executable markup, a URL redirecting to an
internal address, an oversized graph, a forged reviewer role and an action result
that tries to change its run identifier. The expected behavior is rejection or
safe treatment as data, with a recorded reason. These tests protect the authority
boundaries on which the interpretation method relies.

\section{The ordered implementation sequence}

The implementation should proceed in small stages that each produce a reviewable
capability. The first stage freezes contracts and builds deterministic examples.
The next adds persistence and budgets. Only after the controller behaves correctly
with scripted members should it call live model providers. This sequence avoids
confusing stochastic model behavior with defects in state management.

\begin{longtable}{P{13mm}P{49mm}P{85mm}}
\caption{Ordered work packages for the proposed extension.}\label{tab:extension-tasks}\\
\toprule
ID & Deliverable & Acceptance condition\\\midrule\endfirsthead
\toprule ID & Deliverable & Acceptance condition\\\midrule\endhead
E01 & Contract models and reference policy & All valid examples pass; extra fields, incompatible limits and invalid transitions are rejected.\\
E02 & Source inventory and dependency records & Every retained unit has a disposition; missing or mismatched units create issues even under unanimous proposals.\\
E03 & Immutable candidate and assumption store & Revisions retain parents and source hashes; changing a definition invalidates dependent candidates.\\
E04 & Issue identity and lifecycle & Child issues inherit root accounting; terminal unresolved states cannot masquerade as resolutions.\\
E05 & Atomic budget and outbox service & Concurrent reservations never exceed caps; crash recovery preserves issued counts and fencing.\\
E06 & Scripted provider and controller & Omission repair, persistent ambiguity, timeout, no-progress and cancellation traces match the specification.\\
E07 & Structured comparison and witnesses & Different source contexts are detected; supported formal differences replay in Python and Java.\\
E08 & Complete report and review interface & Every issue and mandatory role is represented; dissent and unexplored families remain visible.\\
E09 & Review and release integration & Exact report/candidate/build identity is required; all stale or unresolved release attempts are rejected.\\
E10 & Restricted retrieval and model adapters & No uncounted retries or model-directed tools; malformed and unavailable responses retain distinct outcomes.\\
E11 & Independent reference corpus & Source and scenario judgments are frozen before model outputs; disagreement and adjudication are retained.\\
E12 & Equal-budget evaluation harness & Baselines share source access and cost accounting; uncertainty and abstention metrics are reported together.\\
E13 & Optional argumentation/search extension & Small reference graphs and family-preservation tests pass; incomplete search remains explicit.\\
E14 & Bank integration and shadow operation & Real identity, host mappings, operational dispositions and rollback procedures are approved and exercised.\\
\bottomrule
\end{longtable}

Dependencies matter. E05 depends on the record and issue contracts, while E06
needs E05 and a scripted adapter. E09 depends on complete reports and existing
release identity checks. E10 can be developed behind the adapter interface, but
live evaluation should wait for E11 and E12. E13 is optional until the simpler
scheduler's limitations are measured; it must not delay a clear implementation
of bounded disagreement handling.

Each task should specify files to add or extend, existing interfaces to reuse,
acceptance identifiers and evidence locations. The companion task JSON provides
that structure without changing the existing master plan's completed T00--T22
history. These are new extension tasks, not a retroactive claim that the earlier
MVP failed to implement its agreed scope.

An implementation agent should stop a task for a genuine contract inconsistency,
missing authority or unavailable required infrastructure. It should not stop
merely because a scripted candidate fails: detecting that failure is often the
acceptance criterion. A failed interpretation candidate can trigger the next
planned repair phase while still blocking release. The task record distinguishes
those outcomes.

\section{Acceptance scenarios that exercise the whole cycle}

The first end-to-end fixture uses the explicit type-of-product omission. It
freezes the source packet, emits two candidates, creates a discrepancy, reserves
a repair action, validates the revised formula, reruns affected checks and
produces a report. The report is ready for review only if all other required
coverage and dependencies in that fixture are complete. The real second-circular
packet remains blocked where its dependencies are unresolved.

The second fixture uses the cashback ambiguity with no resolving authority. It
must preserve both supported candidates through the configured rounds and end in
a blocked report. The test checks not only the terminal status but also the
attempt history, retained objections, question for the reviewer and absence of
release authority. This scenario prevents an implementation from treating
exhaustion as permission to choose the highest score.

The third fixture tests unanimous omission. All candidate providers return the
same incomplete rule, while an independent inventory identifies an uncovered
source unit. The controller must open a material issue despite unanimous agreement.
If it does not, the ensemble has implemented voting rather than the proposed
method.

The fourth fixture tests failures under concurrency. Two workers attempt the last
available reservation, one crashes after dispatch, and a stale worker later
returns a result. Exactly one reservation may consume the remaining slot; the
crashed action remains counted; the stale result cannot overwrite a newer state.
The test uses a controlled clock and scripted transport so it is reproducible.
Real distributed-system testing later extends this fixture under the deployment
infrastructure.

The fifth fixture tests invalidation. A reviewer approves a candidate and report,
then a source derivative reveals a previously omitted footnote. The affected
coverage and interpretation records become stale. Release with the earlier
approval must fail until the impact is reviewed and the required checks rerun.
An unchanged Boolean body does not bypass the source-coverage question.

The sixth fixture tests report recovery. A run is cancelled or encounters a
malformed provider response. The service must still produce a complete terminal
report from durable records. If rich rendering fails, a minimal machine-readable
report must remain available. The acceptance condition is not a pretty success
screen; it is a reconstructible account of what happened.

\section{The first vertical slice and its completion rule}

A useful first vertical slice implements E01--E09 with scripted members and the
public second-circular packet. It exercises source inventory, competing candidates,
a repair, a persistent ambiguity, a report and the existing Java path. It does
not need an expensive model integration to test whether the controller's logic
is coherent. This is the smallest product-shaped demonstration of the proposed
method.

Completion requires more than green unit tests. The service must generate an
inspectable report from a fresh run, preserve the exact input and output identities,
show all unresolved issues in the interface, and reject a release attempt for the
blocked example. A separate synthetic fully resolved fixture can demonstrate the
positive review path using local test identities. Its approval status must remain
clearly synthetic.

The implementation note should record the command, environment, source hashes,
policy, result files and any departures from the design. If a departure changes
meaning, budget accounting or approval behavior, it requires a contract revision
and new acceptance cases. Minor internal choices can be made by the agent, but
must not silently alter observable semantics.

Live models then replace scripted proposal adapters without changing the controller
contract. That transition creates a new empirical question: whether the methods
find useful alternatives and omissions on unseen circulars at acceptable cost.
The next chapter supplies the reference process, baselines and uncertainty
analysis needed to answer that question. Engineering completion of the vertical
slice is a prerequisite for that study, not its result.
